\documentclass[11pt,a4paper]{article}

% --- Packages ---
\usepackage[utf8]{inputenc}
\usepackage[margin=1in]{geometry}
\usepackage{amsmath, amsfonts, amssymb}
\usepackage{lmodern}
\usepackage{microtype}
\usepackage{xcolor}
\usepackage{mdframed}
\usepackage{hyperref}

% --- Visual & Typography Setup ---
\definecolor{primary}{RGB}{37, 99, 235}    % Royal Blue
\definecolor{darkgray}{RGB}{51, 65, 85}    % Slate Dark Text
\definecolor{boxbg}{RGB}{248, 250, 252}    % Soft Light Background

\hypersetup{
    colorlinks=true,
    linkcolor=primary,
    urlcolor=primary
}

% Optional: Frame setup for major result equations
\mdfdefinestyle{eqbox}{
    backgroundcolor=boxbg,
    linecolor=primary,
    linewidth=1.5pt,
    topline=false,
    bottomline=false,
    rightline=false,
    innerleftmargin=12pt,
    innerrightmargin=12pt,
    innertopmargin=8pt,
    innerbottommargin=8pt
}

\setlength{\parskip}{0.6em}
\setlength{\parindent}{0pt}

% --- Document ---
\begin{document}

\title{\textbf{\Large \color{primary} The Math Behind This Library}}
\author{}
\date{}
\maketitle

\vspace{-2em}

The neural network is composed of many layers. There are $n$ hidden layers denoted by $[l]$ and for simplicity at this stage we assume there is one output layer $[L]$. Starting from a perceptron, the $i^{th}$ perceptron accepts a vector $a^{[l-1]}$ of size $k$ which represents the output of the \textbf{entire} previous layer. The perceptron/neuron then holds one weight per value of $a^{[l-1]}$ and thus each perceptron holds a vector of weights $W_i^{[l]}$. Each perceptron also holds a bias that we introduce in training the perceptron which is a scalar value of $b_i$. The output $z_i$ for the $i^{th}$ perceptron is the following:

\begin{equation}
z_i^{[l]} = b_i + \sum_{j=1}^{k}{W_{ij}^{[l]} \cdot a_{j}^{[l-1]}}
\end{equation}

To account for non-linear patterns, we pass the $z_i$ value through an activation function $f(z_i)$ defined as:

\begin{equation}
a_i^{[l]} = f(z_i^{[l]})
\end{equation}

To calculate how much our output deviates we defined a loss function $L(a^{[l]})$. Now $a^{[l]}$ is a function of $z^{[l]}$ which is a function of the weights of that layer $W^{[l]}$. Thus Loss can be defined as follows:

\begin{equation}
L(a^{[l]}(z^{[l]}(W^{[l]})))
\end{equation}

To calculate how each layer's output $z^{[l]}$ contributes to the final loss, we need to calculate the error signal for that layer defined as $\delta^{[l]}$:

\begin{equation}
\delta^{[l]} = \frac{\partial{L^{[l]}}}{\partial{z^{[l]}}} = \left(\frac{\partial{L^{[l]}}}{\partial{a^{[l]}}}\right) \left(\frac{\partial{a^{[l]}}}{\partial{z^{[l]}}}\right)
\end{equation}

Given the following equation:

\begin{equation}
a^{[l]} = f(z^{[l]})
\end{equation}

We can calculate $\frac{\partial{a^{[l]}}}{\partial{z^{[l]}}}$ to be:

\begin{equation}
\frac{\partial{a^{[l]}}}{\partial{z^{[l]}}} = {f'(z^{[l]})}
\end{equation}

Substituting that back:

\begin{equation}
\delta^{[l]} = \frac{\partial{L^{[l]}}}{\partial{a^{[l]}}} \odot {f'(z^{[l]})}
\end{equation}

Now how do we calculate $\frac{\partial{L^{[l]}}}{\partial{a^{[l]}}}$? We go back to original perceptron equation where we can see that $z^{[l+1]}$ is a function of $a^{[l]}$:

\begin{equation}
z^{[l+1]} = W^{[l+1]}a^{[l]} + b^{[l+1]}
\end{equation}

\begin{equation}
\frac{\partial{z^{[l+1]}}}{\partial{a^{[l]}}} = W^{[l+1]}
\end{equation}

Now re-defining Loss function to be a function of the next layer $z^{[l+1]}$ which is a function of the present layer's activated output $a^{[l]}$:

\begin{equation}
L(z^{[l+1]}(a^{[l]}))
\end{equation}

\begin{equation}
\frac{\partial{L}}{\partial{a^{[l]}}} = \frac{\partial{L}}{\partial{z^{[l+1]}}} \frac{\partial{z^{[l+1]}}}{\partial{a^{[l]}}}
\end{equation}

\begin{equation}
\implies \frac{\partial{L}}{\partial{a^{[l]}}} = (W^{[l+1]})^T \delta^{[l+1]}
\end{equation}

Therefore, we come with a final expression for the error signal of a layer to be:

\begin{mdframed}[style=eqbox]
\begin{equation}
\delta^{[l]} = \frac{\partial{L}}{\partial{z^{[l]}}} = (W^{[l+1]})^T \delta^{[l+1]} \odot {f'(z^{[l]})}
\end{equation}
\end{mdframed}

Moving back to the loss function, to calculate the weight gradient, we use the following chain:

\begin{equation}
\frac{\partial{L}}{\partial{W^{[l]}}} = \left(\frac{\partial{L^{[l]}}}{\partial{a^{[l]}}}\right) \left(\frac{\partial{a^{[l]}}}{\partial{z^{[l]}}}\right) \left(\frac{\partial{z^{[l]}}}{\partial{W^{[l]}}}\right)
\end{equation}

Substituting $\delta^{[l]}$ into this, we get:

\begin{equation}
\frac{\partial{L}}{\partial{W^{[l]}}} = \delta^{[l]} \left(\frac{\partial{z^{[l]}}}{\partial{W^{[l]}}}\right)
\end{equation}

Given the following:

\begin{equation}
z^{[l]} = W^{[l]} \cdot a^{[l-1]} + b^{[l]}
\end{equation}

\begin{equation}
\frac{\partial{z^{[l]}}}{\partial{W^{[l]}}} = a^{[l-1]}
\end{equation}

Substituting that into our gradient loss equation, we get:

\begin{mdframed}[style=eqbox]
\begin{equation}
\frac{\partial{L}}{\partial{W^{[l]}}} = \delta^{[l]} (a^{[l-1]})^T
\end{equation}
\end{mdframed}

Do not forget the bias!!! We can write the Loss Function as:

\begin{equation}
L(z^{[l]}(b^{[l]}))
\end{equation}

To calculate the bias gradient, calculate the following:

\begin{equation}
\frac{\partial L}{\partial b^{[l]}} = \left(\frac{\partial L}{\partial z^{[l]}}\right) \left(\frac{\partial z^{[l]}}{\partial b^{[l]}}\right)
\end{equation}

\begin{equation}
\frac{\partial L}{\partial b^{[l]}} = \delta^{[l]} \left(\frac{\partial z^{[l]}}{\partial b^{[l]}}\right)
\end{equation}

\begin{equation}
z^{[l]} = W^{[l]} \cdot a^{[l-1]} + b^{[l]}, \quad \frac{\partial z^{[l]}}{\partial b^{[l]}} = 1
\end{equation}

Therefore,

\begin{mdframed}[style=eqbox]
\begin{equation}
\frac{\partial L}{\partial b^{[l]}} = \delta^{[l]}
\end{equation}
\end{mdframed}

\end{document}